\section{Deep-Dream}
\label{sec:deep-dream}

\subsection{Memory representation}

Let $\mathcal D_t$ be the documents processed by time $t$. For a document $D\in\mathcal D_t$, let $\mathcal E(D)$ be the Episodes obtained from its immutable source spans. For an Episode $e_i\in\mathcal E(D_i)$, define one stored observation as
\[
o_i=(D_i,e_i,\tau_i,x_i,\ell_i),
\]
where $\tau_i$ is processing time, $x_i$ is the extracted content and attributes, and $\ell_i$ is the exact source offset in $D_i$. The observation store after $t$ writes is $O_t=\{o_1,\ldots,o_t\}$. The persistent state is
\[
M_t=(F_t,O_t,G_t),
\]
where $F_t$ contains concept families and candidate identities, and $G_t$ contains relation, neighborhood, and version pointers. Every observation carries the document and Episode identifiers $(D_i,e_i)$ that recover its source text.

A write step is represented by
\begin{equation}
\begin{aligned}
F_{t+1}&=\mathsf{Align}(F_t,x_{t+1}),\\
O_{t+1}&=O_t\cup\{o_{t+1}\},\\
G_{t+1}&=\mathsf{Link}(G_t,F_{t+1},o_{t+1}),\\
M_{t+1}&=(F_{t+1},O_{t+1},G_{t+1}).
\end{aligned}
\label{eq:state-transition}
\end{equation}
Here $\mathsf{Align}$ either assigns $x_{t+1}$ to an existing family or keeps a new candidate family; $\mathsf{Link}$ adds or redirects pointers between families, observations, and versions. Neither operation edits the text addressed by $\ell_{t+1}$. Under the representation invariant that version pointers retain the predecessor state, a projection $\pi_t$ can recover $M_t$ from $M_{t+1}$. This is a conditional property of the retained state, not an empirical guarantee that the implementation can recover every historical graph configuration; updates add links without deleting the earlier observation record.

This representation separates finding from reading: resolving a family or relation is an index lookup, while opening its source span is a separate load. A relation edge helps the agent find a connected observation, but the assertion supporting it remains in the source text. Processing time $\tau_i$ is the single state axis of the graph; an event date may appear in $x_i$ and help answer a question, but it does not create a second version clock.

\subsection{Writing memory}

Writing follows the state transition in \cref{eq:state-transition}: documents are split into Episodes with source offsets; entities and relations are extracted from each Episode; candidate identities are compared with existing concept families; and accepted links are written together with their supporting observations. Every link therefore has a path back to a document span.

When the first mention of an entity is ambiguous, Deep-Dream keeps candidate families separate. Later observations and relation neighborhoods can provide enough evidence to merge compatible families or redirect their pointers. The operation changes $F_t$ and $G_t$, not the observations in $O_t$. A temporary split is therefore recoverable, whereas an incorrect early merge can mix two histories under one search handle.

\subsection{Immediate and delayed entity alignment}

The experiments compare two write schedules over the same state representation. \textbf{Immediate alignment (Base)} applies $\mathsf{Align}$ as each window arrives. \textbf{Delayed alignment (Align+)} first appends $o_i$ and candidate links, then resolves compatible families after the surrounding windows have been processed. The difference is the time of the identity decision, not the schema or source record.

Delayed alignment uses more context before making a merge. In our implementation it is bundled with batched or concurrent processing, so the paired measurements in \cref{sec:temporal-benchmark} quantify the deployed bundle rather than delayed alignment in isolation. It can reduce repeated work while preserving the source spans that support each decision.

\subsection{Reading memory}
\label{sec:answer-paths}

Let $\mathcal S_t$ be the set of immutable source spans in $M_t$. For a query $q$, the agent starts with $C_0(q)=\varnothing$ and chooses an action
\[
a_{k+1}=\rho\bigl(q,M_t,C_k(q),u_k\bigr)
\in\bigl(\mathcal S_t\setminus C_k(q)\bigr)\cup\{\mathsf{stop}\},
\]
where $u_k$ is the agent's current uncertainty. If $a_{k+1}$ is a span, it is opened and
\[
C_{k+1}(q)=C_k(q)\cup\{a_{k+1}\};
\]
if $a_{k+1}=\mathsf{stop}$, the chain ends. By induction---each non-stop step adds one member of $\mathcal S_t$ to $C_k$---we have $C_k(q)\subseteq\mathcal S_t$ and $C_k(q)\subseteq C_{k+1}(q)$. The equation tracks source-opening actions; lexical, semantic, concept, and relation lookups update $u_k$ without adding a span to $C_k$. This formalizes the source-opening trace; the policy decides when the evidence is sufficient. The information terms below are identities from the chain rule, not estimates of navigation utility or answer quality.

At query time, the agent begins with the cheapest useful cue---lexical, semantic, concept, or relation---opens the corresponding Episode or source span, and expands the search when the evidence is incomplete: a question about one fact may stop after one passage, while an update or multi-document question may follow several relation and version links.

We evaluate three read paths:
\begin{enumerate}
    \item \textbf{Fixed retrieval} runs the lexical/semantic/graph retriever once and gives the answerer one ranked list of memory records and retrieved chunks. The list cannot be changed or checked against a source file.
    \item \textbf{Memory-tool reading} gives the answering agent read-only search, concept, relation, and chunk tools. It may call them repeatedly and use one result to form the next query, but the answerer receives the records returned by those tools.
    \item \textbf{Source-grounded reading} uses the index to select document groups, opens the original files, reads exact spans, and submits only spans that pass the source gate.
\end{enumerate}

The source check verifies that a cited identifier was returned by a tool and still points to the original text. It establishes provenance of the citation; the agent must still decide whether the passage supports the answer.

\subsection{Why this representation preserves useful detail}
\label{sec:information-view}

Let $Z$ be a fixed target assignment of observations to entities. The dimension of $Z$ is held fixed as observations arrive. Since $\pi_t(M_{t+1})=M_t$, the newer state contains the previous state together with the new observation and any revised links. Therefore
\begin{equation}
\begin{aligned}
H(Z\mid M_t)-H(Z\mid M_{t+1})
&=I(Z;M_{t+1}\mid M_t)\\
&\ge 0.
\end{aligned}
\label{eq:state-uncertainty}
\end{equation}
The first equality holds because $M_t$ is recoverable from $M_{t+1}$---a representation invariant supplied by the retained version links, not by arbitrary graph edges. If the new observation carries no information about $Z$, the difference is zero; when it distinguishes candidate families, the conditional entropy decreases. This is a structural guarantee about retained state, not a measured entropy gain.

For a query $q$, let $A$ be the compact record used to locate candidates, $E_C$ the source spans opened from candidate set $C$, and $Y$ the answer-relevant fact. The chain rule gives
\begin{equation}
I(Y;A,E_C\mid Q=q)=I(Y;A\mid Q=q)+I(Y;E_C\mid A,Q=q).
\label{eq:two-rate-code}
\end{equation}
This is the chain rule. The first term is information used to locate a candidate; the second is the additional information obtained by reading its source. Deep-Dream compresses the search description, not the only copy of the fact.

Let $R_q$ be the event that a relevant source span is in the opened set $C_K(q)$, and let $V_q$ be the event that the produced answer is correct, $V_q=\{\widehat Y=Y\}$. By the law of total probability,
\begin{equation}
\begin{aligned}
\Pr[\widehat Y=Y\mid q]
&=\Pr[R_q\mid q]\Pr[V_q\mid R_q,q]\\
&\quad+\Pr[\neg R_q\mid q]\Pr[V_q\mid\neg R_q,q]\\
&\ge \Pr[R_q\mid q]\Pr[V_q\mid R_q,q].
\end{aligned}
\label{eq:answer-decomposition}
\end{equation}
Concept and relation links primarily affect the first factor; source reading and citation checking affect the second. The inequality is not a performance estimate or a claim that retrieval is necessary for a correct answer; it only decomposes outcomes according to whether a relevant span was opened.
